\chapter{Reference manual}
\label{ch:refman}

\begin{aside}
The rest of this book taught you to think in \maya{}. This
chapter is the dictionary you reach for when you've already
internalised the grammar and just need to remember the
spelling. Every public API surface, every DSL pipe, every
widget, paired with the smallest example that exercises it.
\end{aside}

\section{One header, two namespaces}

Everything below the widget line lives behind one include:

\begin{lstlisting}[language=C++]
#include <maya/maya.hpp>
using namespace maya;
using namespace maya::dsl;
\end{lstlisting}

\noindent Widgets are header-per-widget:

\begin{lstlisting}[language=C++]
#include <maya/widget/badge.hpp>
#include <maya/widget/markdown.hpp>
\end{lstlisting}

\noindent And the low-level pipeline, for the rare case you
need it:

\begin{lstlisting}[language=C++]
#include <maya/internal.hpp>
\end{lstlisting}

\section{The Program concept}

A type \code{P} satisfies \code{Program} when it provides:

\begin{itemize}[leftmargin=*]
  \item nested \code{Model} (any data),
  \item nested \code{Msg} (usually a \code{std::variant}),
  \item \code{static Model init()} \emph{or}
    \code{static std::pair<Model, Cmd<Msg>> init()},
  \item \code{static auto update(Model, Msg)
    -> std::pair<Model, Cmd<Msg>>},
  \item \code{static Element view(const Model\&)},
  \item optionally \code{static Sub<Msg> subscribe(const Model\&)}.
\end{itemize}

\noindent The smallest program that satisfies all six:

\begin{lstlisting}[language=C++]
#include <maya/maya.hpp>
using namespace maya;
using namespace maya::dsl;

struct Hello {
    struct Model {};
    struct Quit {};
    using Msg = std::variant<Quit>;

    static Model init() { return {}; }

    static auto update(Model m, Msg)
        -> std::pair<Model, Cmd<Msg>>
    {
        return {m, Cmd<Msg>::quit()};
    }

    static Element view(const Model&) {
        return text("press q") | Bold;
    }

    static auto subscribe(const Model&) -> Sub<Msg> {
        return key_map<Msg>({ {'q', Quit{}} });
    }
};

static_assert(Program<Hello>);

int main() { run<Hello>({.title = "hello"}); }
\end{lstlisting}

\section{\code{run} and \code{RunConfig}}

\begin{lstlisting}[language=C++]
template <Program P>
void run(RunConfig cfg = {});

struct RunConfig {
    std::string_view title = "";
    bool             mouse = false;
    Mode             mode  = Mode::Fullscreen;
    Theme            theme = theme::dark;
};
\end{lstlisting}

\noindent Designated initialisers; set only what you need:

\begin{lstlisting}[language=C++]
run<App>({});                              // all defaults
run<App>({.title = "editor"});
run<App>({.title = "editor", .mouse = true});
run<App>({.theme = theme::light});
\end{lstlisting}

\section{\code{Cmd<Msg>}: side effects as data}

Returned from \code{update}; performed by the runtime.

\begin{lstlisting}[language=C++]
Cmd<Msg>{}                              // no-op
Cmd<Msg>::quit()                        // exit
Cmd<Msg>::batch({c1, c2})               // multiple
Cmd<Msg>::after(100ms, msg)             // delayed
Cmd<Msg>::task([]{ return Msg{...}; }) // async work
Cmd<Msg>::set_title("editor")
Cmd<Msg>::write_clipboard("text")
\end{lstlisting}

\noindent A working stopwatch using \code{after}:

\begin{lstlisting}[language=C++]
#include <maya/maya.hpp>
#include <chrono>
using namespace maya;
using namespace maya::dsl;
using namespace std::chrono_literals;

struct Stop {
    struct Model { bool flash = false; int n = 0; };
    struct Bump {}; struct Clear {}; struct Quit {};
    using Msg = std::variant<Bump, Clear, Quit>;

    static Model init() { return {}; }

    static auto update(Model m, Msg msg)
        -> std::pair<Model, Cmd<Msg>>
    {
        return std::visit(overload{
            [&](Bump) {
                m.n++; m.flash = true;
                return std::pair{m,
                    Cmd<Msg>::after(300ms, Clear{})};
            },
            [&](Clear) {
                m.flash = false;
                return std::pair{m, Cmd<Msg>{}};
            },
            [](Quit) { return std::pair{Model{}, Cmd<Msg>::quit()}; },
        }, msg);
    }

    static Element view(const Model& m) {
        auto n = text(std::to_string(m.n));
        return v(
            m.flash ? n | Bold | fgc(Color::yellow())
                    : n | Bold,
            text("b bump, q quit") | Dim
        ) | padding(1) | border(BorderStyle::Round);
    }

    static auto subscribe(const Model&) -> Sub<Msg> {
        return key_map<Msg>({ {'b', Bump{}}, {'q', Quit{}} });
    }
};

int main() { run<Stop>({.title = "stop"}); }
\end{lstlisting}

\section{\code{Sub<Msg>}: subscriptions}

\begin{lstlisting}[language=C++]
Sub<Msg>::none()
Sub<Msg>::batch({s1, s2})
Sub<Msg>::on_key([](const KeyEvent& k)
    -> std::optional<Msg> { ... })
Sub<Msg>::on_mouse([](const MouseEvent& m)
    -> std::optional<Msg> { ... })
Sub<Msg>::on_resize([](Size sz)
    -> std::optional<Msg> { ... })
Sub<Msg>::on_paste([](std::string text)
    -> std::optional<Msg> { ... })
Sub<Msg>::every(1s,
    [](auto now) -> Msg { return Tick{now}; })

// Shorthand for the common case
key_map<Msg>({ {'q', Quit{}}, {'+', Inc{}} });
\end{lstlisting}

\noindent A ticking counter that pauses when idle:

\begin{lstlisting}[language=C++]
#include <maya/maya.hpp>
#include <chrono>
using namespace maya;
using namespace maya::dsl;
using namespace std::chrono_literals;

struct Tick {
    struct Model { int n = 0; bool on = true; };
    struct Pulse {}; struct Toggle {}; struct Quit {};
    using Msg = std::variant<Pulse, Toggle, Quit>;

    static Model init() { return {}; }

    static auto update(Model m, Msg msg)
        -> std::pair<Model, Cmd<Msg>>
    {
        return std::visit(overload{
            [&](Pulse)  { m.n++;       return std::pair{m, Cmd<Msg>{}}; },
            [&](Toggle) { m.on = !m.on; return std::pair{m, Cmd<Msg>{}}; },
            [](Quit)    { return std::pair{Model{}, Cmd<Msg>::quit()}; },
        }, msg);
    }

    static Element view(const Model& m) {
        return v(
            text("count: " + std::to_string(m.n)) | Bold,
            text(m.on ? "running" : "paused") | Dim
        ) | padding(1) | border(BorderStyle::Round);
    }

    static auto subscribe(const Model& m) -> Sub<Msg> {
        auto keys = key_map<Msg>({
            {'q', Quit{}}, {' ', Toggle{}},
        });
        if (!m.on) return keys;
        return Sub<Msg>::batch({
            keys,
            Sub<Msg>::every(500ms,
                [](auto) -> Msg { return Pulse{}; }),
        });
    }
};

int main() { run<Tick>({.title = "tick"}); }
\end{lstlisting}

\section{Event helpers}

For use inside \code{on\_key} / \code{on\_mouse} lambdas, or
inside \code{canvas\_run} event handlers:

\begin{lstlisting}[language=C++]
key(ev, 'q')                         // char
key(ev, SpecialKey::Up)              // special
ctrl(ev, 'c')                        // Ctrl+C
alt(ev, 'f')                         // Alt+F
shift(ev, SpecialKey::Tab)           // Shift+Tab
any_key(ev)                          // catch-all
as_key(ev)                           // const KeyEvent*

mouse_clicked(ev, MouseButton::Left)
mouse_released(ev)
mouse_moved(ev)
scrolled_up(ev) / scrolled_down(ev)
mouse_pos(ev)                        // std::optional<MousePos>
as_mouse(ev)                         // const MouseEvent*
\end{lstlisting}

\noindent \code{SpecialKey}: \code{Up}, \code{Down},
\code{Left}, \code{Right}, \code{Home}, \code{End},
\code{PageUp}, \code{PageDown}, \code{Tab}, \code{BackTab},
\code{Backspace}, \code{Delete}, \code{Insert},
\code{Enter}, \code{Escape}, \code{F1}..\code{F12}.

\noindent \code{MouseButton}: \code{Left}, \code{Middle},
\code{Right}, \code{WheelUp}, \code{WheelDown}.

\section{Element DSL: text}

\begin{lstlisting}[language=C++]
// Compile-time, zero allocation
t<"Hello"> | Bold | Fg<255, 100, 80>
"Hello"_t | Bold                     // UDL form

// Runtime
text("hello")
text("hello") | Bold | fgc(Color::red())
text(42)                             // int auto-formats
text(3.14)                           // float too

// Wrap behaviour
text("longish") | clip               // truncate at width
text("longish") | nowrap             // no wrap
\end{lstlisting}

\section{Element DSL: containers}

\begin{lstlisting}[language=C++]
v(child1, child2, child3)            // vertical (column)
h(child1, child2, child3)            // horizontal (row)

// A vector of children works as a child slot:
std::vector<Element> rows = { text("a"), text("b") };
v(t<"Header"> | Bold, rows, t<"Footer"> | Dim)
\end{lstlisting}

\section{Element DSL: compile-time pipes}

For literal arguments — the values are template parameters,
so they resolve at compile time and emit no runtime cost.

\begin{lstlisting}[language=C++]
v(...) | pad<1>                       // all sides
v(...) | pad<1, 2>                    // vertical, horizontal
v(...) | pad<1, 2, 3, 4>              // top, right, bottom, left
v(...) | gap_<1>                      // gap between children
v(...) | grow_<1>                     // flex grow
v(...) | w_<40>                       // fixed width
v(...) | h_<10>                       // fixed height
v(...) | border_<Round>               // border style
v(...) | bcol<50, 54, 62>             // border colour (rgb)
v(...) | bordered<Round, 0x323746>    // border + colour
v(...) | Fg<255, 100, 80>             // foreground rgb
v(...) | Bg<30, 30, 40>               // background rgb
v(...) | fg<0xFF4444>                 // foreground hex
v(...) | bg<0x1A1A2E>                 // background hex
\end{lstlisting}

\section{Element DSL: runtime pipes}

Same syntax with runtime values:

\begin{lstlisting}[language=C++]
v(...) | padding(1)                   // all sides
v(...) | padding(1, 2)                // vertical, horizontal
v(...) | padding(1, 2, 3, 4)          // t, r, b, l
v(...) | margin(1)                    // same variants
v(...) | gap(n)
v(...) | grow(1.0f)
v(...) | width(40) | height(10)
v(...) | border(BorderStyle::Round)
v(...) | bcolor(c)
v(...) | btext("Title")
v(...) | btext("Status",
               BorderTextPos::Bottom,
               BorderTextAlign::Right)
v(...) | fgc(Color::cyan())
v(...) | bgc(Color::rgb(30, 30, 40))
v(...) | align(Align::Center)
v(...) | justify(Justify::SpaceBetween)
v(...) | overflow(Overflow::Hidden)
\end{lstlisting}

\noindent \code{Align}: \code{Start}, \code{Center},
\code{End}, \code{Stretch}, \code{Baseline}, \code{Auto}.

\noindent \code{Justify}: \code{Start}, \code{End},
\code{Center}, \code{SpaceBetween}, \code{SpaceAround},
\code{SpaceEvenly}.

\noindent \code{Overflow}: \code{Visible}, \code{Hidden},
\code{Scroll}.

\section{Element DSL: dynamic helpers}

\begin{lstlisting}[language=C++]
dyn([&] { return text("n = " + std::to_string(n)); })
when(is_loading, spinner, content)    // if/else
when(show_debug, debug_panel)         // if only

map(items, [](auto& s) { return text(s) | Bold; })

text("debug") | visible(debug_mode)

space          // flex-growing spacer
spacer()       // function form
sep            // horizontal separator
vsep           // vertical separator
blank_         // empty line
blank()        // function form
\end{lstlisting}

\section{A complete dashboard layout}

\begin{lstlisting}[language=C++]
#include <maya/maya.hpp>
using namespace maya;
using namespace maya::dsl;

struct Dash {
    struct Model { std::string status = "READY"; };
    struct Quit {};
    using Msg = std::variant<Quit>;

    static Model init() { return {}; }
    static auto update(Model m, Msg)
        -> std::pair<Model, Cmd<Msg>> {
        return {m, Cmd<Msg>::quit()};
    }

    static Element view(const Model& m) {
        return v(
            h(
                text("My App") | Bold,
                space,
                text(m.status) | fgc(Color::green())
            ),
            sep,
            h(
                v(text("Nav 1"), text("Nav 2"))
                    | width(20) | padding(1),
                vsep,
                v(text("Content"))
                    | grow(1) | padding(1)
            ) | grow(1),
            sep,
            text("[q] quit") | Dim
        ) | padding(1) | gap(1);
    }

    static auto subscribe(const Model&) -> Sub<Msg> {
        return key_map<Msg>({ {'q', Quit{}} });
    }
};

int main() { run<Dash>({.title = "dash"}); }
\end{lstlisting}

\section{Style}

\begin{lstlisting}[language=C++]
Style{}                              // empty
Style{}.with_bold()
Style{}.with_dim()
Style{}.with_italic()
Style{}.with_underline()
Style{}.with_strikethrough()
Style{}.with_inverse()
Style{}.with_fg(Color::red())
Style{}.with_bg(Color::rgb(30, 30, 40))

// Compose -- rhs wins on conflicts
auto heading = bold_style | Style{}.with_fg(Color::cyan());

// Predefined: bold_style, dim_style,
//             italic_style, underline_style
\end{lstlisting}

\section{Colour}

\begin{lstlisting}[language=C++]
Color::black()    Color::bright_black()
Color::red()      Color::bright_red()
Color::green()    Color::bright_green()
Color::yellow()   Color::bright_yellow()
Color::blue()     Color::bright_blue()
Color::magenta()  Color::bright_magenta()
Color::cyan()     Color::bright_cyan()
Color::white()    Color::bright_white()

Color::rgb(255, 128, 0)
Color::hex(0xFF8000)            // consteval
Color::indexed(214)             // 256-colour palette
Color::hsl(30.f, 1.f, 0.5f)     // hue, sat, light
\end{lstlisting}

\section{Theme}

\begin{lstlisting}[language=C++]
theme::dark        // default, truecolor
theme::light       // truecolor light
theme::dark_ansi   // 16-colour dark
theme::light_ansi  // 16-colour light

run<App>({.theme = theme::light});

// Custom override at compile time
constexpr auto mine = Theme::derive(
    theme::dark,
    [](Theme& t) { t.primary = Color::hex(0xFF6B6B); }
);
\end{lstlisting}

\noindent Theme fields you can read from \code{ctx.theme}:
\code{primary}, \code{secondary}, \code{accent},
\code{success}, \code{error}, \code{warning}, \code{info},
\code{text}, \code{inverse\_text}, \code{muted},
\code{surface}, \code{background}, \code{border},
\code{diff\_added}, \code{diff\_removed},
\code{diff\_changed}, \code{highlight}, \code{selection},
\code{cursor}, \code{link}, \code{placeholder},
\code{shadow}, \code{overlay}.

\section{Borders}

\begin{lstlisting}[language=C++]
BorderStyle::None
BorderStyle::Single         // standard box drawing
BorderStyle::Double         // double lines
BorderStyle::Round          // rounded corners
BorderStyle::Bold           // heavy lines
BorderStyle::SingleDouble   // mixed
BorderStyle::DoubleSingle   // mixed
BorderStyle::Classic        // ASCII fallback
BorderStyle::Arrow

// Selective sides via the box builder
box().border(BorderStyle::Single)
     .border_sides(BorderSides::all())
     .border_sides(BorderSides::horizontal())
     .border_sides(BorderSides::vertical())
     .border_sides(BorderSides::top())
     // | composes
     .border_sides(BorderSides::top() | BorderSides::left())
\end{lstlisting}

\section{Signals: fine-grained reactivity}

For the parts of an app that don't fit Model-Update-View
cleanly, \maya{} also ships reactive signals.

\begin{lstlisting}[language=C++]
Signal<int> count{0};

count.get()                     // read; auto-tracks
count()                         // shorthand
count.set(42)                   // write
count.update([](int& n) { n *= 2; });
count.version()                 // change counter

auto doubled = computed([&] { return count.get() * 2; });
auto labeled = count.map([](int n) {
    return "value: " + std::to_string(n);
});

Effect e([&] {
    std::println("count = {}", count.get());
});

{
    Batch b;        // RAII batched updates
    x.set(10);
    y.set(20);
}                   // one notify, not two
\end{lstlisting}

\section{\code{print}: one-shot inline render}

\begin{lstlisting}[language=C++]
void print(const Element& root);
void print(const Element& root, int width);
\end{lstlisting}

\noindent A summary card from a build tool:

\begin{lstlisting}[language=C++]
#include <maya/maya.hpp>
using namespace maya;
using namespace maya::dsl;

int main() {
    print(
        v(
            text("build complete") | Bold
                | fgc(Color::green()),
            text("warnings: 0") | Dim,
            text("errors:   0") | Dim
        ) | padding(1) | border(BorderStyle::Round)
    );
}
\end{lstlisting}

\section{\code{live}: inline animation loop}

\begin{lstlisting}[language=C++]
struct LiveConfig {
    int  fps       = 30;
    int  max_width = 0;          // 0 = auto
    bool cursor    = false;
};

template <AnyLiveRenderFn F>
void live(LiveConfig cfg, F&& render_fn);

void quit() noexcept;            // ends live() loop
\end{lstlisting}

\noindent A three-second inline progress bar:

\begin{lstlisting}[language=C++]
#include <maya/maya.hpp>
#include <chrono>
using namespace maya;
using namespace maya::dsl;

int main() {
    auto start = std::chrono::steady_clock::now();
    live({.fps = 30}, [&]() -> Element {
        using namespace std::chrono;
        auto e = duration_cast<milliseconds>(
            steady_clock::now() - start).count();
        if (e >= 3000) quit();
        int filled = static_cast<int>(e * 30 / 3000);
        std::string bar(30, '.');
        for (int i = 0; i < filled && i < 30; ++i) bar[i] = '#';
        return h(
            text(bar) | fgc(Color::green()),
            text(" "),
            text(std::to_string(e / 30) + "%") | Bold
        );
    });
    print(text("done") | Bold | fgc(Color::green()));
}
\end{lstlisting}

\section{Widgets: at a glance}

Every widget satisfies the \code{Node} concept and drops
into the DSL as a child. Pick what you need; include only
the headers you use.

\subsection*{Input}

\code{Input<Cfg>}, \code{TextArea}, \code{Checkbox},
\code{ToggleSwitch}, \code{Radio<N>}, \code{Button},
\code{Select}, \code{Slider}.

\subsection*{Data}

\code{Table}, \code{List<Cfg>}, \code{Tree<T>},
\code{Sparkline}.

\subsection*{Navigation}

\code{Tabs}, \code{Breadcrumb}, \code{Menu},
\code{CommandPalette}.

\subsection*{Display}

\code{Badge}, \code{Callout}, \code{ToastManager},
\code{FileRef}, \code{UserMessage},
\code{AssistantMessage}, \code{ProgressBar},
\code{ActivityBar}, \code{Divider}.

\subsection*{Overlay}

\code{Modal}, \code{Popup}, \code{ScrollableView}.

\subsection*{Visualisation}

\code{BarChart}, \code{LineChart}, \code{Gauge},
\code{Heatmap}, \code{PixelCanvas}.

\subsection*{Specialised}

\code{Markdown}, \code{DiffView}, \code{PlanView},
\code{LogViewer}, \code{Calendar}, \code{KeyHelp},
\code{Disclosure}, \code{ThinkingBlock}.

\subsection*{Tools}

\code{ToolCall}, \code{BashTool}, \code{EditTool},
\code{WriteTool}, \code{ReadTool}, \code{FetchTool},
\code{SearchResult}, \code{AgentTool}, \code{Permission}.

\subsection*{Working widget composition}

\begin{lstlisting}[language=C++]
#include <maya/maya.hpp>
#include <maya/widget/badge.hpp>
#include <maya/widget/callout.hpp>
#include <maya/widget/progress.hpp>
#include <maya/widget/table.hpp>

using namespace maya;
using namespace maya::dsl;

struct App {
    struct Model {}; struct Quit {};
    using Msg = std::variant<Quit>;
    static Model init() { return {}; }
    static auto update(Model m, Msg)
        -> std::pair<Model, Cmd<Msg>> {
        return {m, Cmd<Msg>::quit()};
    }
    static Element view(const Model&) {
        return v(
            Badge::info("status"),
            Callout::success("Build passed",
                             "All 42 tests green"),
            ProgressBar{.value = 0.75f, .label = "Progress"},
            Table({
                .headers = {"Name", "Value"},
                .rows    = {{"fps", "60"}, {"mem", "128MB"}},
            })
        ) | padding(1);
    }
    static auto subscribe(const Model&) -> Sub<Msg> {
        return key_map<Msg>({ {'q', Quit{}} });
    }
};

int main() { run<App>({.title = "widgets"}); }
\end{lstlisting}

\section{Core types}

\begin{lstlisting}[language=C++]
Columns w{80};       Rows h{24};        // strong wrappers
w.value;             // int
w + Columns{5};      // arithmetic preserves type

Size s{80, 24};
s.width.value;  s.height.value;
s.is_zero();    s.area();

Position p{10, 5};
Position::origin();           // {0, 0}

Rect r{{0, 0}, {80, 24}};
r.left(); r.top(); r.right(); r.bottom();
r.contains(p);
r.intersect(other);           // std::optional<Rect>
r.unite(other);

Edges<int> e{1, 2, 3, 4};     // t, r, b, l
e.horizontal();  e.vertical();

Dimension::auto_();
Dimension::fixed(40);         // or px(40)
Dimension::percent(50.0f);

using namespace maya::literals;
50_pct;                       // same as percent(50)
\end{lstlisting}

\section{Errors (low-level only)}

\begin{lstlisting}[language=C++]
enum class ErrorKind {
    Io, ParseError, NotFound, PermissionDenied,
    InvalidInput, Timeout, Unknown,
};

struct Error {
    ErrorKind   kind;
    std::string message;
    static Error from_errno(std::string_view ctx);
};

Result<int> r = ok(42);
Status      s = ok();
Result<int> e = err<int>(Error{ErrorKind::Io, "boom"});

if (r) {
    int v = *r;               // operator* extracts
}
r.value_or(0);
r.error();                    // only if !r

// Macros that early-return on failure
MAYA_TRY(name, expr);         // binds name on success
MAYA_TRY_VOID(expr);          // for Status
\end{lstlisting}

\section{The escape hatch: \code{canvas\_run}}

When the Program shape doesn't fit — direct event polling,
custom timing, low-level canvas access — drop to:

\begin{lstlisting}[language=C++]
#include <maya/internal.hpp>

canvas_run(
    {.title = "demo", .mouse = true},
    [&](const Event& ev) -> bool {
        if (key(ev, 'q')) return false;   // false = quit
        return true;
    },
    [&](const Ctx& ctx) -> Element {
        return text("size: " + std::to_string(
            ctx.size.width.value));
    }
);
\end{lstlisting}

\noindent \code{Ctx} carries \code{size} and \code{theme}.
Everything else is the same DSL you already know.

\section{Putting it together: counter}

The example from the top of \code{API.md}, reproduced here
so the chapter ends as it began — with one complete program
that exercises Program, \code{Cmd}, \code{Sub}, the DSL,
colour, padding, and \code{key\_map}:

\begin{lstlisting}[language=C++]
#include <maya/maya.hpp>
using namespace maya;
using namespace maya::dsl;

struct Counter {
    struct Model { int count = 0; };
    struct Inc {}; struct Dec {}; struct Quit {};
    using Msg = std::variant<Inc, Dec, Quit>;

    static Model init() { return {}; }

    static auto update(Model m, Msg msg)
        -> std::pair<Model, Cmd<Msg>>
    {
        return std::visit(overload{
            [&](Inc)  { return std::pair{Model{m.count + 1}, Cmd<Msg>{}}; },
            [&](Dec)  { return std::pair{Model{m.count - 1}, Cmd<Msg>{}}; },
            [](Quit)  { return std::pair{Model{}, Cmd<Msg>::quit()}; },
        }, msg);
    }

    static Element view(const Model& m) {
        auto c = m.count >= 0 ? Color::green() : Color::red();
        return v(
            text("Count: " + std::to_string(m.count))
                | Bold | fgc(c),
            text("[+/-] change  [q] quit") | Dim
        ) | padding(1);
    }

    static auto subscribe(const Model&) -> Sub<Msg> {
        return key_map<Msg>({
            {'+', Inc{}}, {'-', Dec{}}, {'q', Quit{}},
        });
    }
};

int main() { run<Counter>({.title = "counter"}); }
\end{lstlisting}

\section{Recap}

\begin{recap}
\begin{itemize}[leftmargin=*]
  \item One header (\code{<maya/maya.hpp>}), two namespaces
    (\code{maya}, \code{maya::dsl}).
  \item Program concept: six members; the runtime calls them.
  \item Effects are \code{Cmd<Msg>}; inputs are
    \code{Sub<Msg>}; layout is the DSL.
  \item Compile-time pipes for literals; runtime pipes for
    variables.
  \item Widgets are header-per-widget and drop into the DSL.
  \item For everything else there's \code{canvas\_run} and
    \code{<maya/internal.hpp>}.
\end{itemize}
\end{recap}

\section*{Workshop}

\begin{enumerate}[leftmargin=*]
  \item Build and run every code block in this chapter. They
    are deliberately small.
  \item Pick three widgets you've never used. Compose them
    in a single screen.
  \item Replace one of your existing apps' bespoke layouts
    with the DSL pipes documented above.
\end{enumerate}
